\subsection*{Совершенные графы}
\begin{definition}[Совершенный граф]
Граф \(G\) \emph{совершенный}, если \(\chi(H)=\omega(H)\) для каждого индуцированного
подграфа \(H\le G\). Индуцированный подграф совершенного графа совершенен.
\end{definition}
Всегда \(\chi\ge\omega\). Нечётный цикл \(C_{2k+1}\), \(k\ge2\), не совершенен:
\(\chi=3\), \(\omega=2\). Совершенны: \(K_n\) (\(\chi=\omega=n\)), пустые графы (\(\chi=\omega=1\)),
двудольные графы (\(\chi=\omega=2\) при наличии рёбер; индуцированные подграфы тоже
двудольны), в частности чётные циклы.

\subsection*{Интервальные графы}
\begin{definition}[Интервальный граф]
\(G\) — \emph{интервальный}, если существуют интервалы \(I_1,\ldots,I_n\subset\R\) такие, что
\(v_iv_j\in E\iff I_i\cap I_j\ne\varnothing\). Далее \(I_k=(\ell_k,r_k)\).
\end{definition}
\begin{graphillustration}{Пример интервального графа. Пересекающиеся интервалы дают рёбра; вертикальная прямая пересекает 3 интервала — это $\omega=\chi=3$ (клика $\{2,3,4\}$). Цифра в вершине — номер интервала.}{fig:13-intervals}
\begin{tikzpicture}[font=\small]
\draw[gi interval] (0,2.0) -- (2.0,2.0) node[right] {$I_1$};
\draw[gi interval] (1.2,1.4) -- (3.6,1.4) node[right] {$I_2$};
\draw[gi interval] (2.6,0.8) -- (5.0,0.8) node[right] {$I_3$};
\draw[gi interval] (3.0,0.2) -- (4.2,0.2) node[right] {$I_4$};
\draw[gi interval] (4.6,1.6) -- (6.2,1.6) node[right] {$I_5$};
\draw[->] (-0.3,-0.4) -- (6.9,-0.4) node[right] {$\R$};
\draw[gi removed] (3.3,-0.6) -- (3.3,2.3);
\begin{scope}[shift={(9,1)}]
\node[gi vertex] (1) at (-1.6,0.6) {1};
\node[gi vertex] (2) at (-0.6,0.6) {2};
\node[gi vertex] (3) at (0.6,0.6) {3};
\node[gi vertex] (4) at (0,-0.5) {4};
\node[gi vertex] (5) at (1.6,0.6) {5};
\draw[gi edge] (1) -- (2) -- (3) -- (5); \draw[gi edge] (2) -- (4) -- (3);
\end{scope}
\end{tikzpicture}
\end{graphillustration}

\begin{lemma}[Индуцированные подграфы]\label{lem:int-induced}
Индуцированный подграф интервального графа интервален: достаточно оставить
интервалы его вершин.
\end{lemma}
\begin{lemma}[Свойство Хелли на прямой; \(\omega\)]\label{lem:helly}
Если интервалы попарно пересекаются, то у них есть общая точка. Поэтому
\(\omega(G)\) равно наибольшему числу интервалов, покрывающих одну точку.
\end{lemma}
\begin{proof}
Для \(S\subseteq V\): \(\bigcap_{k\in S}I_k=(\max\ell_k,\min r_k)\). Пусть \(\ell_i=\max\ell_k\),
\(r_j=\min r_k\). Из \(I_i\cap I_j\ne\varnothing\) следует \(\max(\ell_i,\ell_j)<\min(r_i,r_j)\), т.\,е.
\(\ell_i<r_j\), и пересечение непусто. Значит, клика — это множество интервалов с общей
точкой, и наоборот: интервалы через одну точку попарно пересекаются.
\end{proof}
\begin{lemma}[\(\chi\) интервального графа]\label{lem:greedy-int}
\(\chi(G)\) равно наибольшему числу интервалов, покрывающих одну точку, т.\,е. \(\chi=\omega\).
\end{lemma}
\begin{proof}
Жадный алгоритм: перебираем интервалы по возрастанию левых концов и даём
каждому наименьший цвет, не занятый уже покрашенными соседями. Пусть
обрабатывается \(I=(\ell,r)\), \(J=(\ell_J,r_J)\) — ранее покрашенный сосед. Тогда
\(\ell_J\le\ell\) и \(J\cap I\ne\varnothing\) влечёт \(r_J>\ell\). Значит, все такие \(J\) вместе с \(I\)
содержат точки \((\ell,\ell+\varepsilon)\) при малом \(\varepsilon\): это клика, и соседей не больше
\(\omega-1\). Цвет \(\le\omega\), поэтому \(\chi\le\omega\le\chi\).
\end{proof}
\begin{theorem}
Интервальный граф совершенен.
\end{theorem}
\begin{proof}
Каждый индуцированный подграф интервален (лемма~\ref{lem:int-induced}), и для него \(\chi=\omega\) (лемма~\ref{lem:greedy-int}).
\end{proof}
\begin{example}[\(C_n\), \(n\ge4\), не интервален]
Более того, в интервальном графе нет индуцированного цикла длины \(\ge4\). Пусть
\(v_1v_2\ldots v_k\) (\(k\ge4\)) — индуцированный цикл и, не умаляя общности, \(r_1\) —
наименьший правый конец среди его интервалов. Соседи \(v_2,v_k\) пересекают \(I_1\) и
кончаются правее \(r_1\): \(r_2,r_k\ge r_1\), \(\ell_2,\ell_k<r_1\). Тогда точки чуть левее \(r_1\)
лежат и в \(I_2\), и в \(I_k\) — \(v_2v_k\in E\), хорда. Противоречие.
\end{example}
\begin{graphillustration}{Почему цикл длины $\ge4$ не интервален: у $I_1$ наименьший правый конец $r_1$. Оба соседа $v_2,v_k$ пересекают $I_1$ и продолжаются правее $r_1$, поэтому пересекаются между собой около $r_1$ (серая полоса) — появляется хорда $v_2v_k$.}{fig:13-c4}
\begin{tikzpicture}[font=\small]
\fill[black!12] (2.7,-0.2) rectangle (3.0,1.6);
\draw[gi interval] (0.5,1.3) -- (3.0,1.3) node[right] {$I_1$};
\draw[gi interval] (1.5,0.75) -- (4.8,0.75) node[right] {$I_2$};
\draw[gi interval] (2.7,0.2) -- (5.6,0.2) node[right] {$I_k$};
\draw[->] (0,-0.4) -- (6.6,-0.4);
\node[gi note] at (3.0,-0.7) {$r_1$};
\end{tikzpicture}
\end{graphillustration}


\subsection*{Графы сравнимости}
\begin{definition}[Граф сравнимости]
Пусть \(<\) — строгий частичный порядок на \(V\). Граф \(G=(V,\{uv: u<v\text{ или }v<u\})\)
называется \emph{графом сравнимости}. Эквивалентно: рёбра \(G\) можно ориентировать
\emph{транзитивно} (\(a\to b\to c\Rightarrow a\to c\)) и без циклов.
\end{definition}
Двудольные графы — графы сравнимости (все рёбра от одной доли к другой).
Индуцированный подграф графа сравнимости — граф сравнимости (ограничение порядка).
\begin{graphillustration}{Слева: частичный порядок (стрелка от меньшего к большему) и его граф сравнимости. Цвет $c(u)$ — число элементов наибольшей цепи, кончающейся в $u$ (цифра у вершины); цепь $a<c<e$ — клика размера $\omega=3$. Справа: «сеть» — не граф сравнимости.}{fig:13-comparability}
\begin{tikzpicture}[font=\small]
\begin{scope}
\node[gi vertex,label=left:{\scriptsize 1}] (a) at (0,0) {$a$};
\node[gi vertex,label=right:{\scriptsize 1}] (b) at (1.6,0) {$b$};
\node[gi vertex,label=left:{\scriptsize 2}] (c) at (0,1.3) {$c$};
\node[gi vertex,label=right:{\scriptsize 2}] (d) at (1.6,1.3) {$d$};
\node[gi vertex,label=left:{\scriptsize 3}] (e) at (0.8,2.6) {$e$};
\draw[gi arrow] (a) -- (c); \draw[gi arrow] (a) -- (d); \draw[gi arrow] (b) -- (d);
\draw[gi arrow] (c) -- (e); \draw[gi arrow] (d) -- (e);
\draw[gi arrow] (a) to[bend left=35] (e);
\draw[gi arrow] (b) to[bend right=35] (e);
\node[gi note] at (0.8,-1.2) {транзитивная ориентация};
\end{scope}
\begin{scope}[shift={(6,0.3)}]
\node[gi vertex] (a) at (0,0) {$a$};
\node[gi vertex] (b) at (1.8,0) {$c$};
\node[gi vertex] (c) at (0.9,1.5) {$b$};
\node[gi vertex] (a1) at (-1,-0.7) {$a'$};
\node[gi vertex] (b1) at (2.8,-0.7) {$c'$};
\node[gi vertex] (c1) at (0.9,2.7) {$b'$};
\draw[gi arrow] (a) -- (c); \draw[gi arrow] (c) -- (b); \draw[gi arrow] (a) -- (b);
\draw[gi edge] (a) -- (a1); \draw[gi edge] (b) -- (b1); \draw[gi edge] (c) -- (c1);
\node[gi note] at (0.9,-1.5) {ребро $bb'$ не ориентировать};
\end{scope}
\end{tikzpicture}
\end{graphillustration}

\begin{example}[Не граф сравнимости]
Треугольник \(abc\) с висячими вершинами \(a',b',c'\) («сеть»). Транзитивная ориентация
треугольника: \(a\to b\to c\), \(a\to c\) (с точностью до имён). Ребро \(bb'\): если \(b'\to b\),
то нужно \(b'\to c\); если \(b\to b'\), то нужно \(a\to b'\). Ни того, ни другого ребра нет.
\end{example}
\begin{theorem}[Мирский, 1971]
Граф сравнимости совершенен.
\end{theorem}
\begin{proof}
Индуцированные подграфы — тоже графы сравнимости, поэтому достаточно доказать
\(\chi(G)\le\omega(G)\). Клика графа сравнимости — \emph{цепь} (попарно сравнимые элементы),
так что \(\omega\) — длина наибольшей цепи. Положим
\[
c(u)=\max\bigl\{|L|:\ L\text{ — цепь, }u=\max L\bigr\}
\]
(в лекциях: длина наибольшего ориентированного пути, кончающегося в \(u\), плюс один).
Тогда \(c\colon V\to\{1,\ldots,\omega\}\). Если \(u<v\), любую цепь, оканчивающуюся в \(u\),
можно продолжить вершиной \(v\), поэтому \(c(v)\ge c(u)+1\). Значит, смежные вершины
получают разные цвета — правильная раскраска в \(\omega\) цветов.
\end{proof}
\begin{tcolorbox}[colback=cbGrayBack,colframe=cbGray!55]
\textbf{Дилуорс.} Дополнение графа сравнимости тоже совершенно — это следует из
слабой теоремы о совершенных графах (билет~\ref{ticket:14}). На языке порядков:
наибольшая антицепь равна наименьшему числу цепей, покрывающих множество.
\end{tcolorbox}
